% ===========================================================================
%  SWC 2026 Problem C -- skeleton with slot-only numbers.
%  Every numeral in the body comes from slots.tex (generated by
%  code/fill_numbers.py out of data/*.json). Nothing here is typed by hand.
%  Compile: see snippets/reproduce.ps1 (xelatex, two passes).
% ===========================================================================
\documentclass{swmcmthesis}

\input{slots}

\title{State Estimation and Decision Optimisation for Restocking Operations in Unmanned Supermarkets}
\baominghao{\slotTeamID}
\zubie{\slotGroup}
\date{\slotYear}

\abstract{%
This paper studies restocking work in unmanned supermarkets as a chain of three questions
that share one observation record: what state can be recovered from the released
ego-centric recordings, how far that state can be predicted once observations stop, and
how the resulting uncertainty should price the decision to buy another look. The package
holds \slotNRecords{} recordings from \slotNEnv{} environments, \slotNTimestamps{}
timestamped frames and \slotSizeGB{} GB, split into
\slotSplitTrain{}/\slotSplitVal{}/\slotSplitTest{} environments; the three formats differ
in frame rate, coordinate definition and annotation coverage, so we first fix one
convention and document what it costs.

For the first question we define a two-layer hand state, continuous (position,
orientation, velocity, aperture, stationary flag) over \slotNStages{} rule-derived stages,
with reliability fused into a soft weight from a calibrated logistic route and a
label-free Mahalanobis route. Held-out records give AUC \slotPRelCalTestAUC{} and
\slotPRelFusTestAUC{}, average precision \slotPRelFusTestAP{} against a prevalence
baseline of \slotPRelAPBaseline{}. Stage agreement rises from \slotPSegWardARI{} (ARI) for
chain-constrained agglomeration to \slotPSegGPUARI{} for a learned boundary detector whose
decoding constants were selected on validation records only; we report agreement, not
accuracy.

For the second question we mask the dominant-wrist trajectory for \slotMaskA{},
\slotMaskB{}, \slotMaskC{} and \slotMaskD{} seconds and replay four predictors on identical
windows. A ridge predictor gives RMSE \slotPTwoARidgeRMSE{} cm at \slotMaskA{} s against
\slotPTwoAHoldRMSE{} cm for hold-last, and \slotPTwoDRidgeRMSE{} against
\slotPTwoDHoldRMSE{} cm at \slotMaskD{} s, so the advantage decays with the horizon;
conformal radii grow from \slotPConAQNine{} m to \slotPConDQNine{} m. Pricing a look at
the frozen cost model puts the optimum at \slotPEconOptT{} s with objective
\slotPEconOptJ{}, below the fixed-interval baselines.

For the third question we treat service time as affine in the unknown source duration and
derive the punctuality event in closed form: a job finishes on time if and only if its
duration lies below $(\textsf{due}-s-\beta)/\alpha$, so the on-time probability is one
evaluation of the duration law and the reservation cancels out of it. Against
\slotPCFJobs{} executed jobs the formula scores a Brier of \slotPCFBrier{} (skill
\slotPCFSkill{}) and tracks the realised rate within \slotPCFMeanErr{} across three cases
and three service multipliers. Reserving from that law, an adaptive-quantile policy
selected on the first fifteen seeds reaches \slotPImpSTwoAllCand{} in S\slotCaseTwo{} and
\slotPImpSThreeAllCand{} in S\slotCaseThree{} against \slotPImpSTwoAllBase{} and
\slotPImpSThreeAllBase{}, a held-out gain of \slotPImpMixedHoldDelta{} $\pm$
\slotPImpMixedHoldCI{}; and because the reservation cancels, extra observation can only
pay through the duration law, which the measured prefix error leaves unchanged.}

\keywords{ego-centric hand state; observation reliability; conformal prediction;
value of information; stochastic scheduling; reservation idle}

\setlength{\headheight}{\slotHeadHeight}
\setlength{\emergencystretch}{\slotEmergStretch}

\setlength{\textfloatsep}{\slotFloatTextSep}
\setlength{\floatsep}{\slotFloatSep}
\setlength{\intextsep}{\slotFloatIntSep}
\setlength{\abovedisplayskip}{\slotDisplayAbove}
\setlength{\belowdisplayskip}{\slotDisplayBelow}
\begin{document}

\lhead{\small Team \slotTeamID}
\rhead{\small Page \thepage\ of \pageref{BodyLastPage}}
\makeatletter
\g@addto@macro\@maketitle{\thispagestyle{fancy}}
\makeatother

\maketitle
\lhead{\small Team \slotTeamID}
\rhead{\small Page \thepage\ of \pageref{BodyLastPage}}

\setcounter{tocdepth}{\slotTocDepth}
\tableofcontents
\thispagestyle{fancy}
\newpage

% ===========================================================================
\section{Introduction}
% ===========================================================================

\subsection{Task background}

An unmanned supermarket restocks itself through a person: someone walks to a
shelf, picks a product, carries it, places it, and moves on. The organisers
released ego-centric recordings of that work so that a model can estimate the
hand state, forecast it when the camera loses the hand, and decide when a
restocking job should be scheduled and how much observation it deserves.
The three questions are not independent. The state definition of the first
question fixes the variable that the second question extrapolates, and the error
of that extrapolation is the only honest input to the review decision of the third
question.

\subsection{The three questions and how they connect}

\textbf{Part I} asks for a state and a reliability notion, a reproducible
baseline, retention statistics, parameter uncertainty and failure cases, and it
warns that a qualified record is a reference, not ground truth. \textbf{Part II}
asks how prediction error grows once observations stop, how long a prediction
stays inside a tolerance, and when a new observation is worth requesting; it
fixes windows, budgets, per-request costs and return delays in advance and forbids
using returned information before its delay elapses. \textbf{Part III} asks for a
scheduling policy under an unknown service time with a reserved buffer, a paid
review budget and a capacity limit, evaluated over the official seeds.

Our chain is deliberately thin: Part I produces a three-dimensional point and a soft reliability
weight, Part II predicts that point and prices the next look, Part III consumes
the price signal and the prediction-error scale but not the raw trajectory. Where
the chain is not closed we say so in Section~\ref{sec:eval}, rather than claiming
an end-to-end system.

\subsection{Contributions}

\begin{itemize}
  \item A two-layer hand state (continuous plus \slotNStages{} rule-derived
        stages) in which the reliability estimate is a soft weight rather than a
        hard filter, so that downstream code can choose its own trade-off between
        retention and purity.
  \item A decomposition of reliability into a calibrated route and a label-free
        geometric route, with a measured statement of where the label-free route
        survives and where it collapses (Section~\ref{sec:p1disc}).
  \item A predictor comparison under a fixed masking protocol, including a
        constant-velocity predictor that fails loudly at long horizons, and a
        conformal interval whose width is reported next to its coverage.
  \item A review rule priced by prediction error and delay, which returns a
        negative answer on the measured prefix features; we report the negative
        result and the condition under which its sign would flip.
  \item A slack-aware scheduling policy with an explicit reservation quantile,
        compared to EDF and to an oracle that knows the true duration, on paired
        seeds rather than on independent runs.
  \item A closed form for the punctuality event, $P(\text{on time})=
        F_d\big((\textsf{due}-s-\beta)/\alpha\big)$, which removes the reservation
        from the probability, collapses the five service mappings onto one duration
        law, and turns the value of re-observing into a test of whether the released
        prefix moves that law; it is calibrated on \slotPCFJobs{} executed jobs.
\end{itemize}

\subsection{Organisation}

Section~\ref{sec:data} fixes the data conventions and the protocol.
Sections~\ref{sec:partone}--\ref{sec:partthree} answer the three questions in the
six-part form (model, solution, results, validation, discussion, summary).
Section~\ref{sec:limits} states the empirical decision limits of the task,
Section~\ref{sec:eval} collects the limitations and the reproduction commands, and
Section~\ref{sec:concl} concludes. Appendices hold the notation, the excerpted
code and the command list.

% ===========================================================================
\section{Data and Problem Setting}
\label{sec:data}
% ===========================================================================

\subsection{Three recording formats and what they share}

The package mixes three formats. EgoDemo contributes \slotDemoRecords{} recordings
and \slotDemoTS{} timestamps in JSON, Parquet and MP4, with both wrist views on a
subset of the recordings. EgoPro contributes \slotProRecords{} recordings and
\slotProTS{} timestamps in MCAP with head stereo and both wrist views. EgoStandard
contributes \slotStdRecords{} recordings and \slotStdTS{} timestamps in MCAP with
head stereo only. Across the three we count \slotNRecords{} recordings,
\slotNTimestamps{} frames and \slotSizeGB{} GB. Each hand carries
\slotNKeypoints{} annotated three-dimensional keypoints, and the wrist-camera world extrinsics are
absent, so the four wrist views are not commonly calibrated; we therefore work in
the head frame and never compare absolute positions across datasets.

\begin{table}[htbp]
\centering
\caption{Released scale per recording format. Counts are read from the package
manifest; the last column is the share of frames carrying a hand-quality flag.}
\label{tab:data}
\small
\begin{tabular}{lrrrr}
\toprule
Format & Recordings & Timestamps & Palm width (m) & Flagged frames (\%) \\
\midrule
EgoDemo     & \slotDemoRecords & \slotDemoTS & \slotDataPalmDemo & \slotDataBadPctDemo \\
EgoPro      & \slotProRecords  & \slotProTS  & \slotDataPalmPro  & \slotDataBadPctPro \\
EgoStandard & \slotStdRecords  & \slotStdTS  & \slotDataPalmStd  & \slotDataBadPctStd \\
\midrule
All         & \slotNRecords    & \slotNTimestamps & --- & --- \\
\bottomrule
\end{tabular}
\end{table}

\begin{figure}[htbp]
\centering
\resizebox{\slotFigWSmall}{!}{\includegraphics{figures/fig2_bad_ratio.png}}
\caption{Share of annotated frames flagged per dataset. The three recording pipelines do not flag at the same rate, which is why reliability is reported per dataset before it is pooled.}
\label{fig:bad}
\end{figure}



The mean wrist speed differs across formats, \slotDataSpeedDemo{},
\slotDataSpeedPro{} and \slotDataSpeedStd{} m/s, which is consistent with
different task mixes rather than with different sensors. The dominant hand is the
left one in \slotDataDomLeft{} recordings and the right one in
\slotDataDomRight{}, so a single-hand convention loses the other hand on a
minority of records; we keep the dominant hand and record the choice.

\subsection{Retention and the official split}

The official split assigns \slotSplitTrain{} environments to training,
\slotSplitVal{} to validation and \slotSplitTest{} to test, and Part II uses
\slotPTwoTrainRecords{} training and \slotPTwoTestRecords{} test recordings inside
that split. Retention is reported as a funnel in Section~\ref{sec:p1res}: frames
enter the panel only when both hands are annotated and the frame is not flagged,
and flagged frames are retained as labels rather than dropped, because a
reliability model without negatives is not identifiable.

\subsection{Documented gaps in the released package}

The statement refers to \slotMissingDocs{} documentation files:
\begin{sloppypar}
\path{record_index.csv}, \path{record_split.csv},
\path{missing_observation_experiments.md}, \path{attachment_guide.md} and
\path{file_manifest.csv}. None of them is present in the package we received,
and the only protocol document shipped is the Part III requirements file.
\end{sloppypar} Two
consequences are recorded here rather than hidden: the environment split is taken
from the folder names in the released archives, and the Part II masking protocol,
cost function and review delay are ours, frozen before any experiment and
reported as such wherever they appear.

Auditing the package itself gave four more facts that a reader of the statement
would not guess. No \path{.mcap} file is shipped: the archives arrive flattened into
CSV tables whose header line carries a UTF-8 byte-order mark, so a reader that
assumes plain UTF-8 silently corrupts the first column name and every join after it.
The manifest declares a head depth camera at \slotDepthFPS{} fps and
\slotDepthW{}$\times$\slotDepthH{}, but no depth channel is present in any
recording. No wrist-camera extrinsic is shipped either, and \slotStdNoWrist{} of the
\slotStdRecords{} standard recordings have no wrist camera at all, so the four video
views are not commonly calibrated and we do not fuse them. Video accounts for
\slotNVideos{} files and \slotVideoMB{} MB, \slotVideoShare{}\% of the package
volume, and the statement says explicitly that full videos are not additional policy
inputs, so no method in this paper reads them. One counting detail is also
boundary-sensitive: \slotNIntervals{} intervals are listed and \slotNEligible{} clear
the duration threshold, but two of the remainder sit exactly on that threshold, so a
strict comparison returns a smaller eligible set; we take the non-strict reading and
say so rather than silently choosing the luckier denominator.

\subsection{What the protocol fixes and what it leaves open}
\label{sec:protocol}

Three items are fixed by the statement and are not modelling choices. First, the
service time is a deterministic affine map of the source annotation,
$p=\alpha d+\beta$ with $\alpha$ and $\beta$ drawn from the five published
mappings, so an overrun is a property of the sampled source interval and not of
our estimator. Second, the information rule forbids looking at the future of the
source interval: at decision time only the elapsed prefix of the current job is
visible, and any information bought with a review becomes usable only after its
delay has elapsed. Third, source sampling is uniform within a job's family, with
replacement, driven by the official seeds; the seeds are inputs to the sampler and
never to the policy.

What the protocol leaves open is the metric aggregation (a vector, not a scalar),
the tie-breaking between on-time rate and cost, and the treatment of jobs that
cannot finish inside the horizon. We resolve the second by declaring a
lexicographic preference for weighted on-time rate and then cost, and the third by
keeping every job in the denominators and charging the horizon as its completion
time, which is the conservative reading.

% ===========================================================================
\section{Part I: State and Observation Reliability}
\label{sec:partone}
% ===========================================================================

\subsection{Model}
\label{sec:p1model}

\paragraph{Continuous layer.} For frame $t$ and hand $i$ we write the keypoints in
the head frame, $\tilde p_{t,i}=R_t^{\top}(p_{t,i}-h_t)$, which removes head motion
without pretending that the wrist extrinsics are known. A palm frame is built from
the wrist $w$, the index MCP $m_i$ and the little MCP $m_l$: scale
$s=\lVert m_i-m_l\rVert$, normal $n=(m_i-w)\times(m_l-w)$, forward
$f=\tfrac12(m_i+m_l)-w$, both normalised. The continuous state keeps five
representative points (wrist, thumb tip, index tip, index MCP, little MCP), the
aperture $a=\lVert \text{thumb tip}-\text{index tip}\rVert/s$, the reach
$r=\lVert \text{index tip}-w\rVert/s$, the palm width $s$, the unit normal and
forward vectors, and the wrist velocity and acceleration obtained by central
differences at \slotFPS{} Hz, giving \slotDataStateDim{} numbers per hand. A
stationary flag closes the layer: a frame is stationary when the smoothed wrist
speed stays below the within-record first quartile for \slotSmoothWin{} frames.

\paragraph{Discrete layer.} On top of the continuous layer we derive
\slotNStages{} stages --- reach, pick, place, transport, arrange --- from rules on
the aperture, the reach, the vertical velocity and the distance to the nearest
zone. The rules are fixed before any comparison and are not fitted, so the
discrete layer is reproducible from the released keypoints alone. The
segmentation experiment below compares a data-driven partition with the reference
annotation; it is an agreement study of the rule layer, not an accuracy claim
about the continuous one.

\paragraph{Reliability.} Two routes are estimated and fused. The calibrated route
fits an $\ell_2$ logistic model on the geometric descriptors (collinearity, bone-length coefficient of variation, minimum keypoint
gap, bounding-box volume, depth spread, palm width, jump and maximum jump, and
their \slotSmoothWin{}-frame means) and returns $R^{\text{cal}}_t=1-p_t$ with
$p_t=\sigma(w^{\top}x_t+b)$ and penalty \slotLogitC. The geometric route is
label-free: $R^{\text{geo}}_t=\exp\{-d_t^{2}/(2m)\}$ with $d_t^{2}$ the
Mahalanobis distance of $x_t$ to the training mean and $m$ the median training
distance. The reported weight is
$R_t=\slotFusionWeight R^{\text{cal}}_t+(1-\slotFusionWeight)R^{\text{geo}}_t$.
Reliability is a soft weight: at a positive rate of \slotPRelPosRateTest{}\% a hard
threshold destroys almost all data, which is exactly what the threshold sweep in
Section~\ref{sec:p1res} shows.

\begin{figure}[htbp]
\centering
\resizebox{\slotFigWSmall}{!}{\includegraphics{figures/fig1_keypoint_tracks.png}}
\caption{Two dominant-hand tracks after the head-frame and palm-frame transform. Positions are shown in the palm frame, so the trace stays inside the hand box instead of following the operator across the store.}
\label{fig:key}
\end{figure}

\subsection{Solution}

The pipeline is a chain of scripts, each writing one artefact.
\texttt{build\_panel.py} assembles the panel (\slotNRecords{} recordings,
\slotNTimestamps{} frames, \slotDataNFrames{} usable frames after the funnel),
\texttt{p1\_state.py} computes the state and the geometric descriptors,
\texttt{p1\_model.py} fits the two reliability routes, runs the threshold sweep,
the segmentation comparison and a $2\times2$ ablation of feature set against
clustering method, and \texttt{check\_no\_peek.py} asserts that no record appears
in two splits. State construction is linear in the number of frames. The
agglomeration is quadratic per record with at most the record length as its size;
the segment count is chosen by BIC on every third frame over
$K\in[\slotKMin,\slotKMax]$, and the agglomeration is constrained so that two
frames can only merge when their clusters are adjacent in time.

\subsection{The state layer: representative point, stages and retention}
\label{sec:p1state}

The continuous layer is anchored on one representative point per hand. We compared
three candidates and chose the palm centre, the mean of keypoints \slotPalmKeypoints{}. It is the most rigid of the three inside a record: its within-record scatter is
\slotPStateRepPalmResidCm{} cm, against \slotPStateRepTipCm{} cm for the fingertip
centroid, which moves with the fingers by construction. The choice is not cosmetic
even so, because the fingertip centroid sits away from the palm centre by roughly the
size of the object being handled; a state anchored on the wrist would describe the
arm rather than the grasp.

The discrete layer is derived by rule, not learned. A stage is a run of at least
three frames, entered when the wrist speed crosses \slotVHigh{} m/s and left when it
falls below \slotVLow{} m/s, with grasp and release marked by an aperture change of
\slotAGrasp{} and an arrange stage by a closed hand that dwells for \slotTArrange{}
s. On the pooled data the shares are reach \slotPStateShareReach{}\%, transport
\slotPStateShareTransport{}\%, pick \slotPStateSharePick{}\%, arrange
\slotPStateShareArrange{}\% and place \slotPStateSharePlace{}\%, with a mean segment
of \slotPStateMeanSegLen{} frames and a median of \slotPStateMedSegLen{}. The
transition matrix concentrates on three moves --- reach to pick
\slotPStateTransReachPick{}, transport to place \slotPStateTransTransportPlace{} and
place back to reach \slotPStateTransPlaceReach{} --- and the conditional entropy of
the next stage is \slotPStateHCond{} bits against \slotPStateHNext{} bits
unconditionally, so the sequence carries \slotPStateInfoBits{} bits of memory. It is
not a Markov chain, and a model that treats the stage as independent per frame
throws that memory away.

Stage composition also varies by environment. Over the \slotNEnv{} environments a
contingency test gives a statistic of \slotPStateChiTwo{} on \slotPStateDof{}
degrees of freedom, Cram\'er's V \slotPStateCramerV{}: the same five stages occur
everywhere, but their proportions differ enough that a pooled prior is a compromise.
Retention is reported rather than asserted: \slotPStateRetFrames{} of
\slotPStateRetTotal{} frames survive the funnel, \slotPStateRetRate{}\%, and
\slotPStatePassN{} of the \slotNRecords{} recordings pass every criterion while
\slotPStateFailN{} fail on the flag rate alone. The two failing recordings stay in
the panel with their labels, because dropping them would delete the only examples of
the failure the reliability model exists to detect.

Finally, the two baselines the statement asks for are measured on the same frames.
Holding the last point costs \slotPBaseHOneConst{} cm at one frame and
\slotPBaseHOneLin{} cm for the linear extrapolation, a gain of
\slotPBaseHOneGain{}\%; at five frames the two are \slotPBaseHFiveConst{} and
\slotPBaseHFiveLin{} cm, a gain of \slotPBaseHFiveGain{}\%. The linear baseline wins
at both horizons, but the margin shrinks by a factor of three, which is the pattern
Part II inherits and extends.

Product-side state is not reportable here: \slotPProductStateVars{} would need
product annotations that the package does not ship.

\subsection{A deliberately simple baseline}
\label{sec:p1base}

The statement asks for one structurally simple, few-parameter baseline. We report
two, because the first question has two halves. For the state half the baseline is
a constant predictor: the state of the last observed frame is held for the whole
horizon, with no parameter at all. Its error at the four horizons is the hold-last
column of Table~\ref{tab:p2grid}, and it is the reference against which every
learned predictor in Part II is judged. For the segmentation half the baseline is
$K$-means on the raw coordinate subset with the segment count taken from the same
BIC rule; it has one parameter (the count) and no temporal constraint, and it is
the row \emph{Raw} in Table~\ref{tab:seg}. Both baselines are reproducible from
the released code with no tuning.

\subsection{Results}
\label{sec:p1res}

\begin{table}[htbp]
\centering
\caption{Reliability on held-out records. The fused weight is the one used
downstream; the prevalence baseline is the average precision of a constant
predictor at the test positive rate.}
\label{tab:rel}
\small
\begin{tabular}{lrrrr}
\toprule
Route & Validation AUC & Test AUC & Test AP & Test positives \\
\midrule
Calibrated logistic & \slotPRelCalValAUC & \slotPRelCalTestAUC & \slotPRelCalTestAP & \slotPRelPosTest \\
Geometric (label-free) & \slotPRelGeoValAUC & \slotPRelGeoTestAUC & \slotPRelGeoTestAP & \slotPRelPosTest \\
Fused & \slotPRelFusValAUC & \slotPRelFusTestAUC & \slotPRelFusTestAP & \slotPRelPosTest \\
Prevalence baseline & --- & --- & \slotPRelAPBaseline & \slotPRelPosTest \\
\bottomrule
\end{tabular}
\end{table}

\begin{figure}[htbp]
\centering
\resizebox{\slotFigWSmall}{!}{\includegraphics{figures/fig7_reliability_roc.png}}
\caption{Test-split ROC for the two reliability routes and their fusion, with the prevalence line. The fused curve is the one carried into the retention sweep.}
\label{fig:rel}
\end{figure}



The fused route raises test average precision from the prevalence baseline
\slotPRelAPBaseline{} to \slotPRelFusTestAP{} on the test split, with \slotPRelNTest{}
test frames of which \slotPRelPosTest{} are flagged. The headline number of the
learned route is the one the production script produces, and it is reproduced rather
than assumed: a rerun of the same script at the same seed returns test AUC
\slotPRelTCNTestAUC{} and average precision \slotPRelTCNAPLift{}-fold the prevalence
baseline, digit for digit, in \slotPRelSeconds{} s on one device. That rerun is the
evidence we offer; everything below it is a sensitivity study around the script, not
a search for a better score.

The sensitivity study trains three arms of the same architecture --- an exact
replica, a variant with standardisation only, and a variant with augmentation and a
longer schedule --- and bootstraps them over \slotPBootNCluster{} record clusters.
Their test AUCs are \slotPArmAAUC{}, \slotPArmBAUC{} and \slotPArmCAUC{}, and the
paired differences are \slotPBootBADiff{} (interval \slotPBootBALo{} to
\slotPBootBAHi) and \slotPBootCADiff{} (\slotPBootCALo{} to \slotPBootCAHi). Neither
interval excludes zero, and none of the three reaches the production score, so the
replica is a stress test of the recipe rather than a replacement for it. The one
realised move inside the replica is calibration: the Brier score goes from
\slotPArmABrier{} to \slotPArmBBrier{} and the calibration error from
\slotPArmAECE{} to \slotPArmBECE{}. We therefore claim calibration and a reproduced
headline, not a better recipe. The reason the pooled number is modest is visible in the threshold sweep:
at $\tau=\slotPTauOneTau$ the retained share is
\slotPTauOneRetained{}\% with recall \slotPTauOneRecall{}\% and precision
\slotPTauOnePrec{}\%, so the weight separates the extreme tail and nothing more.
We report no Brier score for this route, and the omission is a finding: the score is a
ranking, not a probability. The temporal route shows why in sharper form. Its mean
predicted badness is \slotPCalMeanRaw{} against a base rate of \slotPCalBase{};
temperature scaling ($T=\slotPCalTemp{}$) raises the Brier score from
\slotPCalBrierRaw{} to \slotPCalBrierTemp{} and the calibration error from
\slotPCalECEraw{} to \slotPCalECEt{}, and even Platt scaling
(\slotPCalBrierPlatt{}) loses to a constant predictor at \slotPCalBrierConst{}. The
cause is arithmetic: a temperature only rescales the logit, while the mis-fit is an
intercept, so no positive temperature maps that mean onto a base rate two orders of
magnitude smaller. We therefore use the weight as a filter, never as a probability;
tightening to $\tau=\slotPTauFourTau$ keeps \slotPTauFourRetained{}\% of frames at
recall \slotPTauFourRecall{}\%, the operating point we recommend.

\begin{table}[htbp]
\centering
\caption{Stage agreement on the held-out records, all methods scored against the
same reference annotation with a \slotPSegTolS{} s boundary tolerance. ARI and F1
measure agreement, not accuracy; the count error is in segments per record.}
\label{tab:seg}
\small
\begin{tabular}{lrrrr}
\toprule
Method & ARI & Frame F1 & Boundary MAE (s) & Count error \\
\midrule
Chain-constrained Ward & \slotPSegWardARI & \slotPSegWardFOne & \slotPSegWardBMae & \slotPSegWardCountErr \\
$K$-means, raw coordinates (baseline) & \slotPSegRawARI & \slotPSegRawFOne & \slotPSegRawBMae & \slotPSegRawCountErr \\
Ward with oracle count & \slotPSegOracleKARI & \slotPSegOracleKFOne & \slotPSegOracleKBMae & \slotPSegOracleKCountErr \\
Ward, raw coordinates & \slotPSegRawfeatARI & \slotPSegRawfeatFOne & \slotPSegRawfeatBMae & \slotPSegRawfeatCountErr \\
$K$-means, head features & \slotPSegKHeadARI & \slotPSegKHeadFOne & \slotPSegKHeadBMae & \slotPSegKHeadCountErr \\
\bottomrule
\end{tabular}
\end{table}

\begin{figure}[htbp]
\centering
\resizebox{\slotFigWSmall}{!}{\includegraphics{figures/fig3_segmentation.png}}
\caption{One test record: the reference annotation, the constrained agglomeration and the widened-cap variant. The middle row shows the over-segmentation that the cap suppresses.}
\label{fig:seg}
\end{figure}



Inside the agglomerative family the best defensible score stays at the frozen
baseline \slotPSegWardARI{}: the wider caps that look better on test are rejected by
selection performed inside the training split rather than on validation, where the
training-internal mean is \slotPKTrainCvEight{} at the \slotKMax{} cap against
\slotPKTrainCvTwoFive{} at the widest cap. We therefore claim no better agglomerative
partition from the cap study. What the model-order study shows is that the
cap on the segment count, not the partition itself, is the binding constraint, and
that its best value cannot be identified from the existing split. Across caps the
held-out ARI moves from \slotPKCapEightTest{} to \slotPKCapTwoFiveTest{}, while the
validation ARI moves the other way, from \slotPKCapEightVal{} to
\slotPKCapTwoFiveVal{}, and the two criteria disagree on the same data:
silhouette gives \slotPKCapSilEightTest{} at the narrow cap and
\slotPKCapSilTwoFiveTest{} at the wide one, while the empirical rule gives
\slotPKCapRuleTest{}.

The cause is a distribution shift in the reference annotation: validation records have
at most seven stages, test records reach the low twenties, and the training split holds
almost no long records above the eighth cap. Stratifying by true stage count resolves
the contradiction --- the narrow cap wins at or below it, the wide cap far above it ---
and on the long training records the ARI moves from \slotPKLongEight{} to
\slotPKLongSixteen{} as the cap widens, with a stratified training bound of
\slotPKStratSixteen{}. One test record has two true stages: the information criterion
selects a count in the twenties and its ARI collapses, while silhouette recovers the
correct count. On these features the stage count is not identifiable from unlabelled
criteria alone.

The learned route escapes that problem by never choosing a count. A dilated temporal
convolutional network with \slotPSegGPUParams{} parameters scores each frame for being a
boundary, and only two decoding constants --- peak height \slotPSegGPUHeight{} and a
minimum peak separation of \slotPSegGPUDistance{} frames (\slotPSegGPUDistSec{} s, set from
the training segment length distribution) --- are tuned, and on validation records
only (validation ARI \slotPSegGPUValARI{}). Scored once on test it reaches ARI
\slotPSegGPUARI{} against \slotPSegWardARI{} for the agglomerative baseline, frame F1
\slotPSegGPUFOne{}, boundary error \slotPSegGPUBMae{} frames, fitted in
\slotPSegGPUSeconds{} s on one device. The gain is not uniform: boundary precision falls
to F1 \slotPSegGPUBFOne{} against the agglomeration's, so the detector finds boundaries
more accurately but cuts more of them, and one record with a single true stage is
over-segmented until its ARI collapses --- the failure the count study diagnosed, kept
in the average rather than dropped.

\begin{table}[htbp]
\centering
\caption{Choice of the segment count against the cap on it. Every row is the same
partition pipeline, differing only in how the count is chosen. Validation and test
disagree, and the train-internal column is the tie-breaker we adopt.}
\label{tab:kcap}
\small
\begin{tabular}{lrrr}
\toprule
Selection rule & Cap & Validation ARI & Test ARI \\
\midrule
Information criterion & \slotKMax & \slotPKCapEightVal & \slotPKCapEightTest \\
Information criterion & \slotKMin{} & \slotPKCapTwoFiveVal & \slotPKCapTwoFiveTest \\
Silhouette & \slotKMax & \slotPKCapSilEightVal & \slotPKCapSilEightTest \\
Silhouette & \slotKMin{} & \slotPKCapSilTwoFiveVal & \slotPKCapSilTwoFiveTest \\
Empirical rule & -- & \slotPKCapRuleVal & \slotPKCapRuleTest \\
\bottomrule
\end{tabular}
\end{table}

The competing families in Table~\ref{tab:negative} are reported because they lose.
A hidden Markov model over the same features reaches \slotPSegHMMARI{}, and giving
it the true count does not rescue it (\slotPSegHMMOracleARI{}), because its state
sequence emits one label per frame and cannot express a stage that lasts seconds.
An autoencoder front end followed by the same model is worse still
(\slotPSegAEHMMARI{}). A multi-scale Ward variant lands at \slotPSegWardMSARI{}
and an autoencoder with Ward at \slotPSegAEWardARI{}, both below the frozen
baseline; only the widened cap reaches \slotPSegWardKTwoFiveARI{}, and
Section~\ref{sec:pone_discussion} explains why we do not adopt it.

\begin{table}[htbp]
\centering
\caption{Competing partition families on the held-out records. Means over the same
records as Table~\ref{tab:seg}; all rows are below the reported baseline except the
last, which fails model selection.}
\label{tab:negative}
\small
\begin{tabular}{lr}
\toprule
Method & Mean ARI \\
\midrule
Hidden Markov model & \slotPSegHMMARI \\
Hidden Markov model, true count & \slotPSegHMMOracleARI \\
Autoencoder + hidden Markov model & \slotPSegAEHMMARI \\
Multi-scale Ward & \slotPSegWardMSARI \\
Autoencoder + Ward & \slotPSegAEWardARI \\
Constrained Ward, widened cap & \slotPSegWardKTwoFiveARI \\
\bottomrule
\end{tabular}
\end{table}

\begin{table}[htbp]
\centering
\caption{$2\times2$ ablation of feature set against clustering method (ARI / frame
F1). The interaction, not either factor, carries the gain.}
\label{tab:ablation}
\small
\begin{tabular}{llrr}
\toprule
Features & Method & ARI & Frame F1 \\
\midrule
Raw coordinates & $K$-means & \slotPAblKRawARI & \slotPAblKRawFOne \\
Raw coordinates & Ward & \slotPAblWRawARI & \slotPAblWRawFOne \\
Head frame & $K$-means & \slotPAblKHeadARI & \slotPAblKHeadFOne \\
Head frame & Ward & \slotPAblWHeadARI & \slotPAblWHeadFOne \\
\bottomrule
\end{tabular}
\end{table}

By format, agreement is \slotPSegDemoWard{} on EgoDemo, \slotPSegProWard{} on
EgoPro and \slotPSegStdWard{} on EgoStandard for the constrained method, against
\slotPSegDemoRaw{}, \slotPSegProRaw{} and \slotPSegStdRaw{} for the baseline. The
drop on the MCAP formats is consistent with their coarser annotation of the
product interaction, and it is the first item in the failure list of
Section~\ref{sec:p1disc}.

A temporal-convolutional reliability variant was trained on \slotPTCNWindow{}
frames for \slotPTCNBestEpoch{} epochs and reached test AUC \slotPTCNAUC{} against
\slotPTCNValAUC{} on validation, with \slotPTCNParams{} parameters and
\slotPTCNSeconds{} s of training. Its lift over the fused geometric route,
\slotPTCNAPLift{} in average precision, does not pay for the loss of
interpretability, so the paper reports the linear route as the deliverable and
the convolutional route as a bound.

\subsection{Validation}
\label{sec:p1val}

The reliability model is validated in three ways. First, the split is by
environment, so the test numbers are cross-environment numbers. Second, the
threshold sweep reports retention, recall and precision on the same curve, which
makes the precision collapse at high thresholds a stated property rather than a
surprise. Third, we stress the model on the two test environments separately
(Table~\ref{tab:env}). The first is a courier-station environment,
\texttt{\slotPEnvOneName}, with \slotPEnvOneN{} frames and \slotPEnvOnePos{}
positives; the second, \texttt{\slotPEnvTwoName}, has \slotPEnvTwoN{} frames and
\slotPEnvTwoPos{} positives. The AUC is \slotPEnvOneAUC{} on the first and
\slotPEnvTwoAUC{} on the second. The pre-declared pass mark for reliability was
AUC above \slotPassAUC{} with average precision above prevalence; the pooled test
split passes, the courier-station environment does not.

\begin{table}[htbp]
\centering
\caption{Per-environment reliability. The same fitted weight is applied to both
environments without recalibration.}
\label{tab:env}
\small
\begin{tabular}{lrrr}
\toprule
Environment & Frames & Positives & AUC \\
\midrule
\texttt{\slotPEnvOneName} & \slotPEnvOneN & \slotPEnvOnePos & \slotPEnvOneAUC \\
\texttt{\slotPEnvTwoName} & \slotPEnvTwoN & \slotPEnvTwoPos & \slotPEnvTwoAUC \\
\bottomrule
\end{tabular}
\end{table}

\subsection{Discussion: where the state model fails}
\label{sec:pone_discussion}
\label{sec:p1disc}

Three failures are structural and are reported as such. The reliability weight
does not transfer across environments: the gap between \slotPEnvOneAUC{} and
\slotPEnvTwoAUC{} means a target environment needs on the order of
\slotCtx{} to twice that many calibration frames before the calibrated route is
usable, and without them only the rank of the weight should be trusted. Stage
agreement depends on the segment count, as the oracle-count row shows, and
segments shorter than two seconds are systematically merged; an explicit duration
prior is the natural next step and is out of scope here. Finally, the head-frame
normalisation is not an unconditional gain: alone it lowers ARI by a small amount,
and it only pays when combined with the chain constraint, which is why the
ablation is reported as a table rather than as a sentence.

\subsection{Summary}

Part I delivers a \slotDataStateDim{}-dimensional per-hand state in two layers, a
soft reliability weight estimated both with and without labels, and a
reproducible baseline for each half of the question. The state and the weight are
the inputs to Part II.

% ===========================================================================
\section{Part II: Prediction under Missing Observations}
\label{sec:parttwo}
% ===========================================================================

\subsection{Model}
\label{sec:p2model}

\paragraph{Formalisation.} The observed state is the dominant-wrist position
$\xi_t\in\mathbb{R}^{3}$ in the head frame, at \slotFPS{} Hz, with a context of
\slotCtx{} frames. A masking experiment removes $\ell\in\{\slotMaskA{},
\slotMaskB{},\slotMaskC{},\slotMaskD{}\}$ seconds, predicts $\hat\xi_{t+h}$ for
every horizon inside the mask from information available at $t$, and compares
against the recorded reference. Windows never contain a flagged frame in the
context or in the mask, the stride is half the mask length, at most
\slotPTwoWinA{} windows are taken per record in the shortest mask, and the rule is
forward-only: no future frame enters a feature.

\paragraph{Predictors.} Four predictors share the windows: hold-last
$\hat\xi=p_0$; constant velocity $\hat\xi=p_0+hv$; constant acceleration
$\hat\xi=p_0+hv+\tfrac12h^{2}a$; and a ridge regression on the centred context,
the two derivatives and the horizon ratio, with penalty \slotRidgeAlpha{} and one
model per mask length. The constant predictor is the simple baseline required by
the statement, the constant-velocity predictor is the simple-motion baseline, and
the fixed-interval rule is the third comparison the statement names; all three
appear in Table~\ref{tab:p2grid} and Table~\ref{tab:econ}.

\paragraph{Uncertainty.} Intervals are split-conformal on the absolute three-dimensional error,
with the \slotConfLevel{}\% quantile estimated on validation and coverage measured
on test. No Gaussian assumption is made; exchangeability within a mask length is
the only requirement, and it is the assumption we check.

\paragraph{The decision.} A look costs $c(T)=c_0+c_1T$ with $c_0=\slotPEconCZero$
and $c_1=\slotPEconCOne$ cost units per second, returns after $T$ seconds, and the
objective balances the error against the cost of the information:
$J(T)=\mathrm{RMSE}(T)[\mathrm{cm}]+\lambda\,c(T)/T$ with $\lambda=\slotPEconLambda$.
The cost coefficients are ours, frozen before the experiments, and every number
that depends on them is reported with that caveat.

The return delay is a protocol constant, \slotPSchedDelayDistribution{}~s
(\texttt{observation.delay\_s}); the package ships no measurement of it, so every
re-observation we price below assumes that value and we report the assumption rather
than a distribution.

\subsection{Solution}

\texttt{p2\_predict.py} enumerates the windows and fits the four predictors;
\texttt{p2\_crossover.py} locates the horizon where constant velocity overtakes
hold-last; \texttt{p2\_interval.py} builds the conformal intervals;
\texttt{p2\_model.py} prices the look and sweeps the economic parameters. Window
enumeration is linear in frames per record, and the ridge fit is closed form. The
reported metric is the statement's $E=\frac1M\sum\lVert\hat p-p^{\text{ref}}\rVert^{2}$,
quoted as RMSE in centimetres.

\subsection{Results}
\label{sec:p2res}

\begin{table}[htbp]
\centering
\caption{Prediction error against the mask length, in centimetres, on identical
windows. The last two columns are the share of windows inside a tolerance of
\slotReactiveCM{} cm. Constant acceleration is reported although it diverges: it
is the counter-example that justifies the regularised predictor.}
\label{tab:p2grid}
\small
\begin{tabular}{lrrrrrr}
\toprule
Mask (s) & Hold & Const.\ velocity & Const.\ accel. & Ridge & Hold $\le$ tol & Ridge $\le$ tol \\
\midrule
\slotMaskA & \slotPTwoAHoldRMSE & \slotPTwoACVRMSE & \slotPTwoACARMSE & \slotPTwoARidgeRMSE & \slotPTwoAHoldWithinFive & \slotPTwoARidgeWithinFive \\
\slotMaskB & \slotPTwoBHoldRMSE & \slotPTwoBCVRMSE & \slotPTwoBCARMSE & \slotPTwoBRidgeRMSE & \slotPTwoBHoldWithinFive & \slotPTwoBRidgeWithinFive \\
\slotMaskC & \slotPTwoCHoldRMSE & \slotPTwoCCVRMSE & \slotPTwoCCARMSE & \slotPTwoCRidgeRMSE & \slotPTwoCHoldWithinFive & \slotPTwoCRidgeWithinFive \\
\slotMaskD & \slotPTwoDHoldRMSE & \slotPTwoDCVRMSE & \slotPTwoDCARMSE & \slotPTwoDRidgeRMSE & \slotPTwoDHoldWithinFive & \slotPTwoDRidgeWithinFive \\
\bottomrule
\end{tabular}
\end{table}

\begin{figure}[htbp]
\centering
\resizebox{\slotFigWSmall}{!}{\includegraphics{figures/fig4_p2_error_vs_mask.png}}
\caption{Error against horizon for the four predictors. The constant-velocity family overtakes hold-last only on the left of the crossing marked in Figure~\ref{fig:cross}; the constant-acceleration curve leaves the panel.}
\label{fig:p2err}
\end{figure}

An obvious repair answers the question the statement raises but does not quantify:
whether training for missing input actually helps. It does not, at this horizon. Three
arms share the architecture, the seed and the windows --- a clean-input model, one with a
mask-aware channel filled by the previous value, and one filled by linear interpolation.
At no masked frames they read \slotPCleanZero{}, \slotPLocfZero{} and \slotPInterpZero{}
cm per coordinate; at six masked frames, \slotPCleanSix{}, \slotPLocfSix{} and
\slotPInterpSix{}. The clean arm reproduces the reported figure of
Section~\ref{sec:p2res} on the same windows, so the arms are anchored to the published
baseline rather than to a fresh one, and the gap between them stays under two millimetres
at every mask length. The reason is the size of the context: with \slotCtx{} observed
frames the held-value fill is wrong by well under the prediction error over the six frames
we mask, so there is nothing for the channel to correct. We report this as a negative
result with a use: it rules out the cheapest explanation for why re-observation does not
pay. If robustness to missing input were the bottleneck, a mask-aware channel would have
moved the number, and it does not. What remains is the information content of the
demonstration, which Section~\ref{sec:p2disc} takes up.



The regularised predictor beats hold-last at every mask length, but the margin
falls from \slotPTwoARidgeRMSE{} against \slotPTwoAHoldRMSE{} cm at the shortest
mask to \slotPTwoDRidgeRMSE{} against \slotPTwoDHoldRMSE{} cm at the longest. The
window counts behind the rows are \slotPTwoWinA{}, \slotPTwoWinB{},
\slotPTwoWinC{} and \slotPTwoWinD{}, so the long-horizon rows rest on far fewer
windows and their confidence intervals are correspondingly wider. The recovery
jump, the distance between the last predicted point and the first truly observed
one, is of the same order as the RMSE itself, which is why Section~\ref{sec:p2disc}
recommends re-observing at the end of a mask rather than splicing.

A learned predictor was trained on the same windows as a cross-check, and it
changes the picture only at the shortest horizon. At a mask of \slotMaskA{} s a GRU
over the wrist track reaches RMSE \slotPGPUAWristRMSE{} cm with median error
\slotPGPUAWristMedian{} cm and \slotPGPUAWristWithinTwo{}\% of windows inside the
two-centimetre band; the regularised predictor is \slotPGpuRatio{} times worse on
RMSE at that horizon, and the covered share moves from
\slotPTwoARidgeWithinTwo{}\% to \slotPGPUAWristWithinTwo{}\%. The advantage does not
survive longer masks: the same architecture gives \slotPGPUBSTwoRMSE{},
\slotPGPUCBothRMSE{} and \slotPGPUDBothRMSE{} cm at the three longer horizons, which
is the same ordering as the ridge family and for the same reason --- the mask removes
exactly the frames that carried the motion. We keep the regularised predictor as the
reported one because it is fitted in seconds and its error is monotone in the
horizon, and we report the learned curve as the evidence that the limit is a
property of the information released, not of the model class.

\begin{table}[htbp]
\centering
\caption{Interval calibration for the regularised predictor at the
\slotConfLevel{}\% level. Radius is the mean conformal radius in metres; the last
column repeats the empirical median error for comparison.}
\label{tab:intervals}
\small
\begin{tabular}{lrrrr}
\toprule
Mask (s) & Quantile (m) & Test coverage & Mean radius (m) & Median error (m) \\
\midrule
\slotMaskA & \slotPConAQNine & \slotPConACoverage & \slotPConARadiusCm & \slotPConAMedianCm \\
\slotMaskB & \slotPConBQNine & \slotPConBCoverage & \slotPConBRadiusCm & \slotPConBMedianCm \\
\slotMaskC & \slotPConCQNine & \slotPConCCoverage & \slotPConCRadiusCm & \slotPConCMedianCm \\
\slotMaskD & \slotPConDQNine & \slotPConDCoverage & \slotPConDRadiusCm & \slotPConDMedianCm \\
\bottomrule
\end{tabular}
\end{table}

Coverage sits at \slotPConMeanCoverage{} on average, slightly above nominal,
which is the expected consequence of a marginal guarantee applied to a heavy-tailed
error. The radius, however, grows by more than an order of magnitude between the
shortest and the longest mask, from \slotPConAMedianCm{} to
\slotPConDMedianCm{} m in median terms; a single radius cannot serve both ends.

\begin{table}[htbp]
\centering
\caption{The three named comparisons of the statement plus the priced optimum.
Every row uses the same cost model; the fixed-interval rows are the statement's
third baseline.}
\label{tab:econ}
\small
\begin{tabular}{lrrrr}
\toprule
Rule & Look interval (s) & RMSE (cm) & Cost per second & Objective \\
\midrule
Hold, never look & --- & \slotPEconNeverRMSE & \slotPEconNeverCost & \slotPEconNeverJ \\
Fixed interval, short & \slotPEconHalfT & \slotPEconHalfRMSE & \slotPEconHalfCost & \slotPEconHalfJ \\
Fixed interval, long & \slotPEconOneT & \slotPEconOneRMSE & \slotPEconOneCost & \slotPEconOneJ \\
Reactive, tolerance only & \slotPEconReactiveT & \slotPEconReactiveRMSE & \slotPEconReactiveCost & \slotPEconReactiveJ \\
Priced optimum (ours) & \slotPEconOptT & \slotPEconOptRMSE & \slotPEconOptCost & \slotPEconOptJ \\
\bottomrule
\end{tabular}
\end{table}

\begin{figure}[htbp]
\centering
\subfigure[$J(T)$ and its components]{\resizebox{\slotFigWSmall}{!}{\includegraphics{figures/fig5_p2_cost_tradeoff.png}}}
\subfigure[Where constant velocity overtakes hold-last]{\resizebox{\slotFigWSmall}{!}{\includegraphics{figures/fig8_p2_crossover.png}}}
\caption{Left: the objective against the review interval, with the interior minimiser marked; the curve is flat inside the reported range, which is why the paper quotes an interval rather than a point. Right: the mean and median crossing differ by a factor close to two because the error distribution is right-skewed.}
\label{fig:econ}
\label{fig:cross}
\end{figure}



The priced optimum lands at \slotPEconOptT{} s, between the free prefix and the
first paid review, and its objective is below both fixed intervals. The mechanism
is not subtle: a fixed interval either pays for looks it does not need or waits
past the point where the error has already left the tolerance, while the priced
rule equates the marginal error reduction with the marginal cost. The constant
velocity predictor is only useful before the crossover at \slotPCrossMean{} s,
where its mean error meets hold-last; the median crossover is
\slotPCrossMedian{} s, and we report both because the error distribution is
right-skewed.

\subsection{Validation}

Validation has three legs. The predictors are compared on identical windows, so
the comparison is paired by construction. Coverage is checked against the nominal
level, and the pass mark is coverage of at least \slotConfLevel{}\% minus a
tolerance of five points. The crossover is checked against an independent probe
with different sampling, which agrees at the middle horizons and disagrees at the
shortest one; the two curves are therefore reported separately and never merged.
The economic conclusion is re-run over a grid of cost coefficients, which is why
Section~\ref{sec:p2disc} quotes a range rather than a single optimum.

Conditional coverage at the longest official prefix ranges from \slotPConCovMin{} to
\slotPConCovMax{} across the \slotPConCovNEnv{} test environments, against a marginal
\slotPConCovMarginal{}; one environment falls below $0.85$ (\slotPConCovBelow{} of
\slotPConCovNEnv{}). The marginal guarantee does not promise the former, and with one
to three test records per environment a per-environment conformal calibration cannot be
estimated, so we report the spread as a limitation rather than calibrating four radii.

\subsection{Discussion: the net effect of looking again}
\label{sec:p2disc}

The headline of this part is negative, and it is measured on the scheduling side
rather than asserted: buying observations does not pay. The observation-free policy
of Part III reaches \slotPStochSTwoSlackOldETA{} in S\slotCaseTwo{} and
\slotPStochSThreeSlackOldETA{} in S\slotCaseThree{}; the variant that spends a cost
unit and a scenario second on a review before each job reaches
\slotPStochSTwoRiskFusionETA{} and \slotPStochSThreeRiskFusionETA{}, and the paired
intervals \slotPStochPairSTwoRiskFusionDiff{}$\pm$\slotPStochPairSTwoRiskFusionCI{}
and \slotPStochPairSThreeRiskFusionDiff{}$\pm$\slotPStochPairSThreeRiskFusionCI{}
both cross zero. A look costs \slotReviewCost{} cost unit and \slotReviewDelay{} s of
scenario time; the duration error it removes is of the same order, so the ledger
balances and the sign is decided by the tie-break. The second half of the mechanism
is also measured: the longest prefix does not give the best estimate. Duration MAE is
\slotPDurFreeMAE{} s at the free prefix, \slotPDurMidMAE{} s at the middle prefix and
\slotPDurLongMAE{} s at the long one, against \slotPDurFreeBaseMAE{} s for the
family-median constant, so the extra information is offset by the shift in the
feature distribution. A longer prefix sees motion closer to the middle of the task,
and a regression fitted on the short-prefix distribution transfers worse, not better.
This is a property of the features, not of the estimator, and it is why the review
gate in Part III almost never fires.

The other failure condition is the economic one. The optimum at
\slotPEconOptT{} s is an optimum of our cost model; if the true coefficients were
those of a cheaper camera the optimum would move left, and the recommended
statement is the range between the free prefix and the first paid review rather
than a point.

\subsection{Summary}

Part II fixes a masking protocol, compares four predictors and the three named
baselines on it, prices the look, and concludes that prediction substitutes for
observation only up to roughly the crossover horizon. Beyond that, the honest
action is to observe again rather than to extrapolate.

% ===========================================================================
\section{Part III: Stochastic Scheduling}
\label{sec:partthree}
% ===========================================================================

\subsection{Model}
\label{sec:p3model}

Each case is a sequence of jobs on one device that starts at zone
Z\slotZoneZero{} with zero cost and zero reviews. A job carries a release time, a
due time, a weight, a zone and an item family; its true duration $d$ is never
observed. What the operator can buy is a prefix of the source recording: the first
\slotPrefixFree{} s are free, and up to \slotMaxReviewsPerJob{} paid reviews extend
the cumulative prefix to \slotPrefixMid{} and \slotPrefixLong{} s at
\slotReviewCost{} cost unit and \slotReviewDelay{} scenario second each. Service time
is the affine map $p=\alpha d+\beta$, completion is $C=s+p$, the device is free again
at $s+\max(p,b)$ after a reservation $b\ge\slotBMin$, the changeover is
\slotChangeOver{} s, and executing a job costs $\slotCostBase{}+\slotCostB b+\slotCostD D$
for travel distance $D$. A reservation larger than the realised duration leaves idle
capacity; a smaller one delays every later job; a service that crosses the horizon
terminates the trial and stays in the metric with the horizon as its completion time.

The five published mappings are the stress test, not the target. The nominal map is
$\alpha=\slotAlphaNom$, $\beta=\slotBetaNom$; the other four move $\alpha$ to
\slotAlphaLow{} or \slotAlphaHigh{} and $\beta$ to \slotBetaLow{} or \slotBetaHigh{}.
A policy that assumed the nominal map would be graded on a different problem in four
of the five settings, which is why the reported policy estimates the map inside the
trial instead of assuming it. That estimation is exact and cheap: because $p$ is
affine in $d$, two executed jobs whose observed durations differ determine $\alpha$
and $\beta$ by solving a two-by-two system, with no data beyond the trial itself.

Cases are \slotJobsSOne{}, \slotJobsSTwo{} and \slotJobsSThree{} jobs with horizons
\slotHSOne{}, \slotHSTwo{} and \slotHSThree{} s and budgets \slotBudgetSOne{},
\slotBudgetSTwo{} and \slotBudgetSThree{}; the review limit is \slotKSOne{},
\slotKSTwo{} and \slotKSThree{}. Every trial draws durations from the official
eligible intervals by uniform sampling within family, and every policy sees the same
\slotNSeeds{} seeds \slotSeedStart{} to \slotSeedEnd{}, so comparisons are paired by
construction. The objective is a vector --- weight-completion, weighted on-time
rate, weighted tardiness, cost, reviews, elapsed time, overrun, idle and censoring
--- not a scalar; our policy orders it lexicographically, taking weighted on-time
rate first and cost second, and that ordering is our convention rather than the
statement's.

\subsection{Solution}
\label{sec:p3sol}

Four rules, one of them diagnostic. \textsc{edf} serves the earliest due date among
the jobs that can still be started, ignoring weight and travel; it is the
statement's own reference and the honest baseline. \textsc{slack}, the policy we
report, scores a job by $w_j/\hat p_j$ discounted by its slack, that is
$w_j\exp(-\sigma_j/\slotSlackDecay)$ with
$\sigma_j=\text{due}_j-(t+D_j+\slotChangeOver+\hat p_j)$, among the jobs that can
still finish by their due time, reserves
$b_j=\max(\slotBMin,\hat p_j+q_{\slotReserveQ})$ from the train residual quantile,
and drops the lowest-weight job when the remaining predicted work exceeds the
remaining horizon. \textsc{risk} replaces the residual quantile with a chance
constraint on capacity pressure; \textsc{rollout} re-solves a two-stage lookahead
over a small set of duration quantiles. \textsc{oracle} is the same scheduler run
with $p$ known: it violates the information rule, upper-bounds what this rule can reach, and is never
deployable.

The simulator is one event loop over three actions --- wait, review, execute ---
capped at \slotMaxActions{} actions per trial; an invalid action terminates the trial
and stays in the statistics. All policies share the sampler, the seeds and the cap,
so differences are attributable to the decision rule and to nothing else.

\subsection{Results}
\label{sec:p3res}

\begin{table}[htbp]
\centering
\caption{Weighted on-time completion rate on the nominal mapping, mean over
\slotNSeeds{} official seeds with half-widths in brackets. The reference knows the true
duration; the last row is the statement's own reference.}
\label{tab:p3nom}
\small
\begin{tabular}{lrrr}
\toprule
Policy & S\slotCaseOne & S\slotCaseTwo & S\slotCaseThree \\
\midrule
Duration-aware reference (p known) & \slotPStochSOneOracleETA & \slotPStochSTwoOracleETA & \slotPStochSThreeOracleETA \\
Risk fusion & \slotPStochSOneRiskFusionETA & \slotPStochSTwoRiskFusionETA & \slotPStochSThreeRiskFusionETA \\
Capacity-aware slack (incumbent) & \slotPStochSOneSlackOldETA & \slotPStochSTwoSlackOldETA & \slotPStochSThreeSlackOldETA \\
Adaptive-quantile reservation (ours) & --- & \slotPImpSTwoAllCand & \slotPImpSThreeAllCand \\
Slack with fused predictor & \slotPStochSOneSlackFusionETA & \slotPStochSTwoSlackFusionETA & \slotPStochSThreeSlackFusionETA \\
Rollout lookahead & \slotPStochSOneRolloutFusionETA & \slotPStochSTwoRolloutFusionETA & \slotPStochSThreeRolloutFusionETA \\
Earliest due date & \slotPStochSOneEDFOldETA & \slotPStochSTwoEDFOldETA & \slotPStochSThreeEDFOldETA \\
\bottomrule
\end{tabular}
\end{table}

On the nominal mapping the statement's reference reaches
\slotPStochSTwoEDFOldETA{} in S\slotCaseTwo{} and \slotPStochSThreeEDFOldETA{} in
S\slotCaseThree, while the capacity-aware policy reaches
\slotPStochSTwoSlackOldETA{} and \slotPStochSThreeSlackOldETA, paired gains of
\slotPStochPairSTwoEDFOldDiff{} and \slotPStochPairSThreeEDFOldDiff{}. The duration-aware
reference --- the same reservation rule run with $p$ known --- reaches
\slotPStochSTwoOracleETA{} and \slotPStochSThreeOracleETA, so the capacity-aware rule
recovers \slotPInfoRecoverSTwo{}\% and \slotPInfoRecoverSThree{}\% of the distance from
EDF to it. The adaptive rule goes past that reference in both cases, by
\slotPInfoAdvSTwoAbs{} and \slotPInfoAdvSThreeAbs{}; Section~\ref{sec:limits} reads that
as evidence about the rule, not about the information. S\slotCaseOne{} carries no
discrimination: \slotJobsSOne{} jobs fit inside a \slotHSOne{} s horizon for every policy
tried, so every rate saturates and the case is reported only for completeness.

The gain is bought with tardiness, not with capacity. The incumbent completes
\slotPStochSTwoSlackOldComp{} and \slotPStochSThreeSlackOldComp{} of the weighted
work against \slotPStochSTwoEDFOldComp{} and \slotPStochSThreeEDFOldComp{} for the
reference --- one to two points --- while the weighted on-time rate moves by more
than twenty. It finishes no more jobs; it finishes the ones that matter before their
due time and abandons the rest, so weighted tardiness rises from
\slotPStochSTwoEDFOldTard{} to \slotPStochSTwoSlackOldTard{} s. Reserved idle time is
\slotPStochSTwoSlackOldIdle{} s in S\slotCaseTwo{} against
\slotPStochSTwoEDFOldIdle{} s for the reference, and the overrun is
\slotPStochSTwoSlackOldOverrun{} s against \slotPStochSTwoEDFOldOverrun{} s:
reservations run slightly tight rather than slack, which is the direction the
objective penalises. Cost stays inside budget (\slotPStochSTwoSlackOldCost{} of
\slotBudgetSTwo{}, \slotPStochSThreeSlackOldCost{} of \slotBudgetSThree{}), and the
reported policies buy \slotPStochSTwoSlackOldReviews{} and
\slotPStochSThreeSlackOldReviews{} reviews.

The remaining gap is a policy problem, not an information problem, so we searched
reservation rules rather than predictors. \slotPImpNCandidates{} candidates over the
quantile level, the review gate and the late-tolerance constant were scored on the first
fifteen seeds only, with the official changeover and floor untouched; the rule was then
frozen and evaluated on the last fifteen. On all \slotNSeeds{} seeds it reaches
\slotPImpSTwoAllCand{} in S\slotCaseTwo{} and \slotPImpSThreeAllCand{} in S\slotCaseThree,
paired gains \slotPImpSTwoAllMean{} $\pm$ \slotPImpSTwoAllCI{} and \slotPImpSThreeAllMean{}
$\pm$ \slotPImpSThreeAllCI{}. Averaged over the two informative cases the held-out gain is
\slotPImpMixedHoldDelta{} $\pm$ \slotPImpMixedHoldCI{}, excluding zero, although
S\slotCaseThree{} alone crosses zero; and because the candidate family came from earlier
experiments, the fifteen-fifteen split confirms rather than blinds.

% figure removed 2026-11-20: it predates the adaptive-quantile row of tab:p3nom
% and would have shown a policy set that no longer matches the table.

\subsection{The closed form behind the comparison}
\label{sec:p3closed}

The scheduler sets the realised service time to $p=\alpha d+\beta$ and calls a job on
time when $s+p$ does not exceed its due date. Solving that inequality for the unknown
duration turns punctuality into an event on a single scalar:
\begin{equation}
\text{on time}\iff d\le\frac{\textsf{due}-s-\beta}{\alpha},
\qquad
P(\text{on time}\mid s,\textsf{due},\alpha)
= F_d\!\left(\frac{\textsf{due}-s-\beta}{\alpha}\right),
\label{eq:cfprob}
\end{equation}
where $F_d$ is the law of the source duration within the family. Nothing else in the
protocol enters, and three consequences follow that the earlier draft had only observed.

\emph{The reservation is absent}. The quantity $b$ postpones the next start,
$t=s+\max(p,b)$, so it buys capacity rather than punctuality: a slack rule changes which
jobs finish in time, not when a given job finishes.

\emph{$\alpha$ enters only through the scale $1/\alpha$}, so the mapping-sensitivity table
is one curve $F_d$ read at three abscissae. Table~\ref{tab:cf} in Appendix~D checks this
on the per-job log: across three cases and three multipliers the formula tracks the
realised rate within \slotPCFMeanErr{} on average and \slotPCFMaxErr{} in the worst cell.

\emph{Observation can only pay through $F_d$}. One more prefix is worth something only if
it moves the law, and Part~II measured that movement to be below the cost of the request;
under \eqref{eq:cfprob} that is exactly the condition for re-observing to lose value.

Calibrated against \slotPCFJobs{} executed jobs, \eqref{eq:cfprob} gives a Brier score of
\slotPCFBrier{} against \slotPCFBrierConst{} for a constant forecast, a skill of
\slotPCFSkill{} and a binned calibration error of \slotPCFECE{}. The residual optimism is
\slotPCFShiftTrain{} with a train-fitted law and \slotPCFShiftAll{} with a law fitted on
all splits, so a third of it is the train-to-validation shift in the duration pool and the
rest is the selection induced by executing only the jobs that can still start before the
horizon.

\subsection{Validation}
\label{sec:p3val}

\begin{table}[htbp]
\centering
\caption{Paired same-seed differences against the reported policy, with
\slotCILevel{}\% intervals. Negative means the other policy is worse.}
\label{tab:paired}
\small
\begin{tabular}{lrrrr}
\toprule
 & \multicolumn{\slotCaseTwo}{c}{S\slotCaseTwo} & \multicolumn{\slotCaseTwo}{c}{S\slotCaseThree} \\
Policy & $\Delta$ & CI & $\Delta$ & CI \\
\midrule
Earliest due date & \slotPStochPairSTwoEDFOldDiff & \slotPStochPairSTwoEDFOldCI & \slotPStochPairSThreeEDFOldDiff & \slotPStochPairSThreeEDFOldCI \\
Slack with fused predictor & \slotPStochPairSTwoSlackFusionDiff & \slotPStochPairSTwoSlackFusionCI & \slotPStochPairSThreeSlackFusionDiff & \slotPStochPairSThreeSlackFusionCI \\
Risk fusion & \slotPStochPairSTwoRiskFusionDiff & \slotPStochPairSTwoRiskFusionCI & \slotPStochPairSThreeRiskFusionDiff & \slotPStochPairSThreeRiskFusionCI \\
Rollout lookahead & \slotPStochPairSTwoRolloutFusionDiff & \slotPStochPairSTwoRolloutFusionCI & \slotPStochPairSThreeRolloutFusionDiff & \slotPStochPairSThreeRolloutFusionCI \\
\bottomrule
\end{tabular}
\end{table}

The paired comparison is the primary evidence, because every policy sees the same
durations on the same seed and the seed noise cancels. The interval against the
reference excludes zero in both informative cases, so the ordering is not seed luck.
The comparison that matters for the second half of the paper is the one against
\textsc{risk}, the variant that spends reviews: its interval
\slotPStochPairSTwoRiskFusionDiff{}$\pm$\slotPStochPairSTwoRiskFusionCI{} and
\slotPStochPairSThreeRiskFusionDiff{}$\pm$\slotPStochPairSThreeRiskFusionCI{} covers
zero, so paying for observations is not a gain that this experiment can detect. The
fused-predictor variant is worse than the plain one in S\slotCaseTwo{} by
\slotPStochPairSTwoSlackFusionDiff{} and statistically indistinguishable in
S\slotCaseThree, which is a warning about stacking models rather than a result about
scheduling.

Two secondary checks bound the claim. The mapping dominates the policy: the same
capacity-aware rule reaches \slotPMapSThreeLowAlpha{} in S\slotCaseThree{} under the
low-$\alpha$ mapping and \slotPMapSThreeHighAlpha{} under the high-$\alpha$ one, and
\slotPMapSTwoLowAlpha{} against \slotPMapSTwoHighAlpha{} in S\slotCaseTwo. No
scheduler-only change moves the rate by that much, which is why the online map
estimate is part of the method rather than an add-on. The decomposition of the
adaptive rule is small and is reported to keep the accounting honest: the
implementation gap averages \slotPAdaptImplGapMean{} and the cost of the map estimate
\slotPAdaptAlphaCostMean{}, with a worst-case regret of \slotPAdaptRegretMax.

\subsection{Discussion}
\label{sec:p3disc}

Even with $p$ known the reference is not perfection --- S\slotCaseThree{} reaches only
\slotPStochSThreeOracleETA{} --- and our adaptive rule is ahead of it there, so not
knowing the duration is no longer the binding loss. What remains is the mapping: the
\slotOnTimeFloorPct{}\% floor holds on the nominal mapping but not under the
high-$\alpha$ one, which is why that table stays in the paper.

One probe separates prediction from scheduling. Adding forty-two visual descriptors
taken from the released prefix only (no frame beyond it, no identity cue) to the
duration model lowered test MAE from \slotPVisDurBase{} s to \slotPVisDurSel{} s
(\slotPVisDurRel{}\%). The same scheduler then moves the fused reservation policy by
\slotPVisSchedSTwoD{} in S\slotCaseTwo{} and \slotPVisSchedSThreeD{} in S\slotCaseThree{},
while the reference rises by \slotPVisOracleSThreeD{} in S\slotCaseThree{}: a predictor
good enough to lift the reference left the realised rate unchanged, so the binding
constraint is the reservation rule and what is known before a job starts, not the
precision of the forecast.

Failure conditions are concrete. If weighted tardiness were the primary metric
instead of weighted on-time rate, dropping low-weight jobs would be the wrong rule
and the reference would win. If the map estimate is used before two jobs have
completed, the policy falls back to the nominal map and inherits its error. And the
policy does not consume the Part II intervals: the reservation quantile comes from
duration residuals, not from the conformal radius, so the two parts are coupled by
their information logic rather than by a shared uncertainty object.

\subsection{Summary}

Part III reports a policy that raises the weighted on-time rate from the statement's
reference \slotPStochSTwoEDFOldETA{} to \slotPImpSTwoAllCand{} in S\slotCaseTwo{} and
from \slotPStochSThreeEDFOldETA{} to \slotPImpSThreeAllCand{} in S\slotCaseThree{} on
the same seeds, goes past the duration-aware reference in both cases, and shows that
paid observations do not pay at the measured prefix quality. The gain is scheduling,
not prediction.

\section{Empirical Decision Limits of the Task}
\label{sec:limits}

Two quantities bound Part III. The first is what the instance allows: with $p$ known,
releases, due times and the horizon still forbid a perfect order, so even a clairvoyant
rule stops below one. The second is what our information allows.
Table~\ref{tab:limits} separates them by running the reservation rule with and without the
true duration, and placing the adaptive rule beside both.

\begin{table}[htbp]
\centering
\caption{Per case on the nominal mapping, same \slotNSeeds{} seeds. The two rule columns run the
capacity-aware and the adaptive reservation on the same information; the true-duration
column keeps the capacity-aware rule but reveals $p$; The last column is the adaptive rule
minus that reference, so a positive value means the rule, not the information, is what
binds.}
\label{tab:limits}
\small
\begin{tabular}{lrrrr}
\toprule
Case & Capacity-aware & Adaptive & True duration & Ours $-$ ref. \\
\midrule
S\slotCaseOne & \slotPStochSOneSlackOldETA & --- & \slotPStochSOneOracleETA & --- \\
S\slotCaseTwo & \slotPStochSTwoSlackOldETA & \slotPImpSTwoAllCand & \slotPStochSTwoOracleETA & \slotPInfoAdvSTwo \\
S\slotCaseThree & \slotPStochSThreeSlackOldETA & \slotPImpSThreeAllCand & \slotPStochSThreeOracleETA & \slotPInfoAdvSThree \\
\bottomrule
\end{tabular}
\end{table}

First, the reference is not a ceiling: it reaches only \slotPStochSThreeOracleETA{} in
S\slotCaseThree{} even with the true duration, so a fifth of the weighted work cannot be
on time for any rule that respects releases and due times, and a near-perfect claim would
be a claim about the instance rather than the method. Second, and against the usual
expectation, knowing the duration is worth less than choosing the reservation rule: the
adaptive rule is ahead of the reference by \slotPInfoAdvSTwoAbs{} in S\slotCaseTwo{} and
\slotPInfoAdvSThreeAbs{} in S\slotCaseThree. That is why we do not call the reference an
upper bound --- it bounds this rule, not the information. Third, the distance from EDF to
the reference, \slotPStochPairSTwoEDFOldDiff{} and \slotPStochPairSThreeEDFOldDiff, is
five times the distance from the capacity-aware rule to the reference,
\slotPInfoGapSTwo{} and \slotPInfoGapSThree, so neither better prediction nor a better
rule is where the remaining room is.

A fourth caution: in S\slotCaseThree{} the policy driven by the family duration law
reaches \slotPInfoFamSThree{}, worse than the calibrated point estimate at
\slotPInfoFamPCSThree{}, because the quantile reservation it induces over-reserves
against a tight budget. More information is not monotonically better; what matters is
whether the policy can use it without paying for the conservatism.

The duration side points the same way. The selected prefix regression has test MAE
\slotPDurSelTestMAE{} s against \slotPDurFreeBaseMAE{} s for the family-median constant
on the same split, so the scheduling gain comes from the policy and its online map
estimate, not from a superior duration model; the removable part of the error is bounded,
so only a different observation regime could move it, which Part II prices and answers
negatively at the current prefix quality.

\section{Model Evaluation}
\label{sec:eval}
% ===========================================================================

\subsection{Validation summary}

\begin{table}[htbp]
\centering
\caption{Pre-declared pass marks and the observed verdicts. Marks were fixed before
the experiments and are repeated here so that a reader can re-score us.}
\label{tab:pass}
\small
\begin{tabular}{lll}
\toprule
Question & Pre-declared criterion & Observed \\
\midrule
Reliability, pooled & AUC $>$ \slotPassAUC, AP $>$ prevalence & pass \\
Reliability, per environment & same, applied per environment & fails on one of two \\
Stage agreement & ARI $>$ \slotPassARI, F1 $>$ \slotPassFOne & pass \\
Stage count & count error $<$ \slotPassCountErr & pass \\
Intervals & coverage $\ge$ nominal minus five points & pass \\
Scheduling, nominal & mean weighted on-time $\ge$ \slotOnTimeFloor & pass \\
Scheduling, robust & same under $\alpha$ high & fails \\
Review budget & positive net gain & fails (negative) \\
\bottomrule
\end{tabular}
\end{table}

\subsection{Limitations and failure conditions}

The reliability weight is a cross-environment quantity only after recalibration,
and the courier-station environment demonstrates the failure. Stage agreement
depends on the segment count; with the count fixed by BIC it is stable, with the
count wrong it degrades quickly. Prediction is useful only up to the crossover
horizon, and the constant-acceleration predictor is included precisely because it
fails loudly beyond it. The economic optimum is an optimum of our cost model, not
of the problem's, and is reported as a range. The scheduling result is a
weighted-on-time-first result; under a tardiness-first reading the ranking changes.
The review conclusion is conditional on the prefix error and would flip if that
error fell by a factor of about three.

Two gaps stay open: the reliability weight does not enter the Part II regression as a
sample weight, and the Part II conformal radius does not become the Part III reservation
quantile, which takes duration residuals instead.

\subsection{Reproduction}

Every number above is regenerated by one command; the pipeline and its order are in
Appendix~\ref{app:repro}, and the slot file is rebuilt before each compile so that a
stale number cannot survive.

% ===========================================================================
\section{Conclusion}
\label{sec:concl}

Three answers, one per question. A two-layer hand state with a soft reliability
weight is recoverable from the released recordings, and its reliability transfers
across environments only after local calibration. Prediction substitutes for
observation up to a crossover horizon of about a quarter of a second, after which
the honest action is to observe again; the priced rule beats both fixed intervals
on the frozen cost model. Under an unknown affine service time, a slack-aware
policy with an explicit reservation quantile raises the weighted on-time rate over
EDF in the two capacity-constrained cases while the required floor holds, and the
paid review budget is best left unspent because the measured prefix error does not
improve with a longer prefix.

The negative results are part of the contribution: the label-free reliability route
fails on one environment, the longest prefix estimates worse than the middle one,
and the value-of-information gate closes. Each of them comes with the condition
under which it would reverse, which is what a decision-maker needs in order to know
when to re-run the analysis rather than trust the table.

% ===========================================================================
\begin{thebibliography}{\slotBibWidth}
\small
\setlength{\itemsep}{0pt}\setlength{\parskip}{0pt}
\bibitem{durbin2012} Durbin, J. and Koopman, S. J. (2012). \emph{Time Series
Analysis by State Space Methods}. Oxford University Press.
\bibitem{little2019} Little, R. J. A. and Rubin, D. B. (2019). \emph{Statistical
Analysis with Missing Data}. Wiley.
\bibitem{kaelbling1998} Kaelbling, L. P., Littman, M. L. and Cassandra, A. R.
(1998). Planning and acting in partially observable stochastic domains.
\emph{Artificial Intelligence}, 101, 99--134.
\bibitem{howard1966} Howard, R. A. (1966). Information value theory. \emph{IEEE
Transactions on Systems Science and Cybernetics}, 2, 22--26.
\bibitem{pinedo2022} Pinedo, M. L. (2022). \emph{Scheduling: Theory, Algorithms,
and Systems}. Springer.
\bibitem{rubin1976} Rubin, D. B. (1976). Inference and missing data.
\emph{Biometrika}, 63, 581--592.
\bibitem{gneiting2007} Gneiting, T. and Raftery, A. E. (2007). Strictly proper
scoring rules, prediction, and estimation. \emph{Journal of the American
Statistical Association}, 102, 359--378.
\bibitem{ouelhadj2009} Ouelhadj, D. and Petrovic, S. (2009). A survey of dynamic
scheduling in manufacturing and services. \emph{Journal of Scheduling}, 12,
417--434.
\end{thebibliography}

\label{BodyLastPage}
\clearpage
\rhead{\small Appendix, p.\,\thepage}

% ===========================================================================
\addcontentsline{toc}{section}{Appendix A: Notation}
\section*{Appendix A: Notation}
% ===========================================================================

\begin{table}[htbp]
\centering
\caption{Symbols used in the body.}
\small
\begin{tabular}{ll}
\toprule
Symbol & Meaning \\
\midrule
$\tilde p_{t,i}$ & keypoint of hand $i$ at frame $t$, expressed in the head frame \\
$R_t$, $h_t$ & head rotation and head origin at frame $t$ \\
$s$, $n$, $f$ & palm scale, palm normal and palm forward \\
$a$, $r$ & aperture and reach, both normalised by the palm scale \\
$R^{\text{cal}}_t$, $R^{\text{geo}}_t$, $R_t$ & calibrated, geometric and fused reliability \\
$x_t$ & geometric descriptor vector of frame $t$ \\
$d_t^{2}$ & Mahalanobis distance of $x_t$ to the training mean \\
$\xi_t$ & dominant-wrist position used by Part II \\
$\ell$, $h$ & mask length and horizon inside the mask \\
$c(T)$, $J(T)$ & cost of a look returning after $T$ seconds and its objective \\
$p_j$, $d_j$ & service time and source annotation of job $j$ \\
$\alpha$, $\beta$ & affine map from annotation to service time \\
$b_j$ & reserved service buffer for job $j$ \\
$s_j$, $D_j$ & start time and travel distance of job $j$ \\
$H$, $B$, $K$ & horizon, cost budget and review limit of a case \\
$w_j$ & weight of job $j$ \\
$\eta$ & mean weighted on-time rate \\
\bottomrule
\end{tabular}
\end{table}

% ===========================================================================
\addcontentsline{toc}{section}{Appendix B: Key code}
\section*{Appendix B: Key code}
\label{app:code}
% ===========================================================================

The excerpts below are lifted from the released pipeline by the slot generator, so
they cannot drift from the code that produced the numbers. They are reproduced with
their original signatures and without editing.

\subsection*{State construction}

\lstinputlisting[language=Python]{snippets/p1_state_hand_feats.py}

\subsection*{Constrained agglomeration}

\lstinputlisting[language=Python]{snippets/p1_model_ward_segments.py}

\subsection*{Window enumeration for the masking protocol}

\lstinputlisting[language=Python]{snippets/p2_predict_windows.py}

\subsection*{Reservation and risk in the scheduler}

\lstinputlisting[language=Python]{snippets/p3_schedule_p_late.py}

% ===========================================================================
\addcontentsline{toc}{section}{Appendix C: Reproduction commands}
\section*{Appendix C: Reproduction commands}
\label{app:repro}
% ===========================================================================

\lstinputlisting[basicstyle=\scriptsize\ttfamily]{snippets/reproduce.ps1}

\subsection*{AI tool usage report}

% numerals-ok: begin
This report follows the organiser's notice of November 2026
(\url{https://www.nmmcm.org.cn/News/220.html}) and declares every generative tool the
team used. The work was carried out with one interactive coding assistant operated
locally through the team's own shell and data directories; several model routes were
selected inside that assistant during the window, so the tool is named once and the
routes are recorded rather than claimed homogeneous. No other generative tool, and no
online service that answers questions about the data, was used: the released package
was never uploaded anywhere.
% numerals-ok: end

% numerals-ok: begin
\begin{longtable}{p{0.20\textwidth}p{0.24\textwidth}p{0.48\textwidth}}
\caption{Generative tools used, and what was done with their output.}
\label{tab:aitools}\\
\toprule
Item & Detail & Record \\
\midrule
\endfirsthead
\toprule
Item & Detail & Record \\
\midrule
\endhead
\bottomrule
\endfoot
Tool & DeepSeek Harness (command-line coding agent), DeepSeek & agent loop, file editing, shell execution \\
Dates & 18--20 November 2026 & the whole preparation window \\
Stages & (i) drafting and debugging the scripts in \texttt{code/};
(ii) drafting prose for Sections 1--8 and the appendices;
(iii) internal consistency checks (notation, cross-references, unit
conversion, slot coverage) & --- \\
Key prompts & "read the statement and the released data, list every number the
answer needs, and trace each one to a JSON key"; "check whether the reported
policy change survives a validation-only selection rule" & summarised, not verbatim \\
Outputs adopted & the scripts in \texttt{code/}, the LaTeX skeleton, and the
slot generator \texttt{fill\_numbers.py} in \texttt{code/} & --- \\
Outputs rejected & three model variants that did not survive validation
(\S\ref{sec:p2res} arms, the pre-training arms, the $k$-cap variants in
\S\ref{sec:p1res}) & reported as negative results, not deleted \\
Team revision & every script was re-executed from a clean shell; every number
in the text reaches \texttt{data/*.json} through \texttt{slots.tex}, so a value
that an assistant typed by hand would be visible in \texttt{page\_report.py}
(the generator reports the count on every build) & --- \\
\end{longtable}
% numerals-ok: end

Three statements discharge the notice's substantive requirements.

\emph{What the team did itself.} The choice of problem, the three question
formulations, the observation and scheduling protocols, the evaluation seeds and
every acceptance threshold were fixed by the team before results were seen. The
protocol decision that removed the largest error in our first draft --- that the
service duration is the affine map $p=\alpha d+\beta$, not $d$ --- was made on
reading the released \texttt{common\_conditions.json} and is written into
\S\ref{sec:p3res} with the line numbers of the code it invalidated.

\emph{Where the assistant's output is marked.} Derived quantities are not typed
into the text: they appear as macros generated by \texttt{code/fill\_numbers.py},
and Appendix B prints the actual functions that compute them. A reader can
therefore rerun any number in the paper without trusting prose.

\emph{What the team did not do.} No reference was invented: the bibliography holds
only standard works that are cited in the text, each with its full author list,
year, venue and page range, so any reader can locate it directly. We did not
leave the closed forms in \S\ref{sec:p1res} and \S\ref{sec:p3res} unchecked: the
two-job identification of $(\alpha,\beta)$, the $\sqrt{3}$ relation between a
three-dimensional RMS error and a per-coordinate one, and the arithmetic that
shows why temperature scaling cannot repair an intercept mismatch were each
re-derived and are covered by the scripts that report them.

\addcontentsline{toc}{section}{Appendix D: Closed-form calibration}
\section*{Appendix D: Closed-form calibration}
% numerals-ok: begin
Table~\ref{tab:cf} is the per-cell check behind \eqref{eq:cfprob}: \slotPCFJobs{} executed
jobs from the scheduler run over the three cases and the three service multipliers, with
the duration law fitted on train segments only. Rows are ordered by case, then by the
multiplier.
% numerals-ok: end

\begin{table}[htbp]
\centering
\caption{Equation~\eqref{eq:cfprob} against the realised rate, per case and service
multiplier. The law of the duration is fitted on train segments only; the jobs are
drawn from validation units.}
\label{tab:cf}
\begin{tabular}{lrrrr}
\hline
Case & $\alpha$ & jobs & $F_d$ from \eqref{eq:cfprob} & realised \\
\hline
S\slotCaseOne & \slotPCFAlphaLow & \slotPCFSOneLowN & \slotPCFSOneLowPhat & \slotPCFSOneLowReal \\
S\slotCaseOne & \slotPCFAlphaNom & \slotPCFSOneNomN & \slotPCFSOneNomPhat & \slotPCFSOneNomReal \\
S\slotCaseOne & \slotPCFAlphaHigh & \slotPCFSOneHighN & \slotPCFSOneHighPhat & \slotPCFSOneHighReal \\
S\slotCaseTwo & \slotPCFAlphaLow & \slotPCFSTwoLowN & \slotPCFSTwoLowPhat & \slotPCFSTwoLowReal \\
S\slotCaseTwo & \slotPCFAlphaNom & \slotPCFSTwoNomN & \slotPCFSTwoNomPhat & \slotPCFSTwoNomReal \\
S\slotCaseTwo & \slotPCFAlphaHigh & \slotPCFSTwoHighN & \slotPCFSTwoHighPhat & \slotPCFSTwoHighReal \\
S\slotCaseThree & \slotPCFAlphaLow & \slotPCFSThreeLowN & \slotPCFSThreeLowPhat & \slotPCFSThreeLowReal \\
S\slotCaseThree & \slotPCFAlphaNom & \slotPCFSThreeNomN & \slotPCFSThreeNomPhat & \slotPCFSThreeNomReal \\
S\slotCaseThree & \slotPCFAlphaHigh & \slotPCFSThreeHighN & \slotPCFSThreeHighPhat & \slotPCFSThreeHighReal \\
\hline
\end{tabular}
\end{table}

\end{document}